\section{Introduction}

Transposable elements (TEs) are a major component of the human genome and are studied through several complementary representations. Repeat-family resources describe classification and consensus sequence models, genome annotations record individual or representative loci, expression datasets measure transcription in specific samples, and the biomedical literature reports associations with genes, diseases and other entities. These representations are not interchangeable. A family consensus is not the sequence of every genomic copy, an observed genomic locus does not establish activity, and a reported association does not by itself establish causality. The interpretation of TE evidence therefore depends on both the source and the unit of observation.

Established resources provide substantial depth within particular evidence types. Repbase and Dfam support repeat-family classification and reference sequence models \citep{Bao2015Repbase,Kojima2018HumanRepbase,Wheeler2013Dfam}. HERVd, dbRIP, euL1db and TranspoGene organize endogenous retrovirus sequences, retrotransposon insertion polymorphisms, sample-specific L1HS insertions or gene-context annotations, respectively \citep{Paces2002HERVd,Wang2006dbRIP,Mir2015euL1db,Levy2008TranspoGene}. More recent resources have focused on disease or expression evidence. HervD Atlas curates HERV-disease associations and exposes them through an interactive knowledge graph, CancerHERVdb concentrates on HERV expression in cancer, and TE-SCALE provides single-cell TE expression and TE-gene co-expression across human cancers \citep{Li2024HervDAtlas,Stricker2023CancerHERVdb,Meng2026TESCALE}. These specialist resources answer important questions at resolutions that a general integration resource should not claim to replace.

The remaining practical problem is cross-layer navigation. A researcher starting from a TE name may need to establish its classification, inspect a reference sequence or genomic record, compare expression contexts, identify correlated genes and open the papers supporting a disease or gene relation. When these steps require separate resources, names and evidence boundaries must be reconciled repeatedly. Conversely, merging heterogeneous records into a single undifferentiated graph risks making weak or indirect evidence appear stronger than it is. The useful integration task is therefore not only to connect records, but also to retain their provenance and interpretation limits.

Here we describe TE-KG, a human TE-centred resource that connects literature-derived biomedical relations, taxonomy, representative genomic and sequence records, expression profiles and context-specific co-expression networks. TE-KG provides graph, path, table, visualization, download and natural-language access routes to these layers. The central design principle is that integration should make evidence jointly accessible without erasing what each source can support. We report the current database content, data-processing and runtime architecture, and the principal scientific access workflows. We also state the limitations that constrain biological interpretation and the additional release information required before public deployment.

